\clearpage\section{Inferring the heavy mass from squeezed triangles}
\subsection{Frequency, phase and the smooth background}
The spectator outside the horizon has $\Sigma\propto x^{3/2\pm i\mu}$ with $\mu^2=M^2+g^2-9/4$ from Eq.~(115). Long and short modes cross the horizon at times separated by $\ln(k_S/k_L)$. The propagation between these scales supplies the characteristic factor $r^{3/2\pm i\mu}$. Using $r^{i\mu}=\exp(i\mu\ln r)$, a real combination of the two branches takes the form
\[
 T(r)=r^{3/2}\,[A\cos(\mu\ln r)+B\sin(\mu\ln r)].
\]
The two coefficients absorb an unknown amplitude and phase. Since $A\cos u+B\sin u=D\cos(u-\phi)$, $D=\sqrt{A^2+B^2}$ and $\phi=\operatorname{atan2}(B,A)$. A single triangle measures a mixture of phase, background and amplitude. The mass enters through the change of phase across several $r$ values.

The fitted shape is $F_{\rm fit}=b_0+b_2r^2+T(r)$, optionally adding $b_4r^4$ to test background sensitivity. At fixed trial $\mu$, make a design matrix with one row per triangle:
\[
 X_i=\big(1,\ r_i^2,\ r_i^{3/2}\cos(\mu\ln r_i),\ r_i^{3/2}\sin(\mu\ln r_i)\big).
\]
Let $y_i=F(r_i)$ and let $s_i$ be its assigned error scale. Divide row $X_i$ and datum $y_i$ by $s_i$, then solve the linear least-squares problem for $c=(b_0,b_2,A,B)^T$. This gives
\[
 \chi^2_{\rm prof}(\mu)=\min_c\sum_i[(y_i-X_i c)/s_i]^2.
\]
Scan $0.1\le\mu\le4$ with spacing $0.01$, refine each candidate minimum, and take the smallest residual. This profiles over the unknown amplitude, phase and background at each frequency. Once $g$ is specified, take the positive root
\begin{equation}
 M=\frac mH=\sqrt{\hat\mu^2+9/4-g^2}.
\end{equation}
Frequency alone constrains the combination $M^2+g^2$. The conversion here uses the input $g=0.1$. At input $M=2$ the target is $\mu=\sqrt{4+0.01-2.25}=1.32665$.

\subsection{Known-input trials}
Use 50 log-spaced ratios in $[0.001,0.2]$. The injected coefficients are $(b_0,b_2,A,B)=(0,0.2,0.1,0.05)$. Add independent Gaussian noise with standard deviation $1\%$ or $5\%$ of the RMS of $T$ on this grid. Fit 100 realisations at each noise level. Coverage counts the fraction of trials whose $\Delta\chi^2\le3.841459$ profile set contains the injected frequency.
\begin{center}\begin{tabular}{lrrrr}\toprule Noise & Mean $\hat\mu-\mu$ & Frequency RMSE & Coverage & Held-out false alarms\\\midrule
1\% & 0.00243 & 0.04687 & 93/100 & 2/100\\
5\% & 0.02179 & 0.20504 & 94/100 & 4/100\\\bottomrule\end{tabular}\end{center}
\clearpage
The null trials contain only the analytic background and noise. On 100 calibration trials per noise level, use the 95th percentile of $\chi^2_{\rm background}-\min_\mu\chi^2_{\rm prof}$ as a detection threshold. Apply it to 100 new null trials at each noise level. This checks how often searching over frequency fits an apparent clock to noise alone.
\fig{mock_mass_recovery}{.91}{Recovered frequencies for known-input mocks (left) and example frequency profiles (right). Increasing the noise broadens the distribution of fitted frequencies.}
\subsection{Physical fit and sensitivity to analysis choices}
Compute $F$ for the same 50 ratios with the coupled modes. At each point calculate its regulator error directly. Interpolate the remaining pilot errors in $\ln r$, add them to the regulator error, and use a floor of $10^{-14}$ for the fit weight. These scales quantify numerical sensitivity in the deterministic calculation.

Using the full window $[0.001,0.2]$ and background $b_0+b_2r^2$ gives $\hat\mu=1.31610$, $D=1.3890\times10^{-7}$ and
\[
 M=\sqrt{1.31610^2+2.25-0.01}=1.9930,\qquad
 N_{\rm cycles}=\frac{\hat\mu\ln(0.2/0.001)}{2\pi}=1.110 .
\]
The largest numerical error relative to the fitted envelope is
$\max_i[\delta F_i/(D r_i^{3/2})]=0.008524$. Thus the envelope exceeds the numerical error scale by a factor of at least $117$ on this grid. The comparison is to the envelope because the oscillatory function crosses zero. The narrow frequency range spanned by roughly one cycle makes it important to vary the smooth terms.
\clearpage
\subsection{What bounds the mass estimate?}
\fig{physical_clock_fit}{.87}{Physical shape scaled by $r^{3/2}$ (left) and the full-window frequency profile (right). The reference line marks the predicted frequency. The profile uses numerical weights.}
Repeat the fit with two lower endpoints, three upper endpoints and both analytic backgrounds. Each entry below is the best frequency on the grid points inside the indicated window.
\begin{center}\begin{tabular}{rrrr}\toprule $r_{\min}$ & $r_{\max}$ & $b_0+b_2r^2$ & $b_0+b_2r^2+b_4r^4$\\\midrule
0.001 & 0.05 & 1.32317 & 1.32544\\
0.003 & 0.05 & 1.31765 & 1.32421\\
0.001 & 0.10 & 1.32041 & 1.32445\\
0.003 & 0.10 & 1.31637 & 1.32142\\
0.001 & 0.20 & 1.31610 & 1.32196\\
0.003 & 0.20 & 1.35637 & 1.31787\\\bottomrule\end{tabular}\end{center}
The largest shift occurs in the last row. Removing the most squeezed points shortens the phase baseline, while retaining $r$ as large as $0.2$ gives the analytic terms more leverage. Adding $r^4$ then moves the optimum from $1.35637$ to $1.31787$. For the full window, adding $r^4$ reduces the numerically weighted residual sum from $181.09$ to $6.45$. A precise integral can still leave a poor fit if the smooth component is too restrictive.

Across the twelve fits, $1.31610\leq\hat\mu\leq1.35637$. Because $M(\mu)$ increases for positive $\mu$, the smallest and largest masses follow directly from the endpoint frequencies:
\begin{equation}
 \sqrt{1.31610^2+2.24}=1.9930,\qquad
 \sqrt{1.35637^2+2.24}=2.0198.
\end{equation}
The input $M=2$ lies inside this range. Its endpoints describe variation under the tested background and window choices. Only the injected-noise mocks above have the noise model needed for a statistical coverage statement. The physical calculation supplies neither survey data nor a likelihood for measurements.
\source{9}{Mock trials, physical triangles, regulated values and the twelve fitted models}
